% =====================================================================
%  APÉNDICE — PRÁCTICA INTEGRADORA DE UNIDAD II
%  CRUD con patrón DAO «sin framework» en JSP/JSTL/EL, Django y Laravel
%  (Reutiliza el preámbulo del libro. Requiere \lstdefinestyle{python}.)
% =====================================================================

\appendix

% --- Formato de capítulo para apéndices: «Apéndice X» en vez de «Semana X»
\titleformat{\chapter}[display]
  {\normalfont\Huge\bfseries\color{uteqverde}}
  {\filleft\Large Apéndice \thechapter}
  {1ex}{\titlerule\vspace{1ex}\filright}

\chapter{Práctica integradora: CRUD con patrón DAO sin framework en JSP/JSTL/EL, Django y Laravel}
\label{ap:crud-dao}

\epigraph{«El mismo problema, resuelto tres veces, enseña más que tres
problemas resueltos una sola vez: lo que permanece es el patrón; lo que
cambia es el lenguaje.»}{--- \textit{Principio pedagógico de la
asignatura}}

\section*{Resultado de aprendizaje}
\addcontentsline{toc}{section}{Resultado de aprendizaje}
Al concluir esta práctica, el estudiante construye una misma aplicación
\emph{CRUD} (crear, leer, actualizar y eliminar) sobre la tabla
institucional \texttt{estudiantes} en tres pilas tecnológicas del lado
del servidor ---Java con JSP/JSTL/EL, Django y Laravel---, aplicando en
las tres el patrón \emph{Data Access Object} (DAO) y el acceso a datos
con SQL parametrizado escrito a mano, sin delegar la persistencia en un
ORM. Con ello consolida la separación de responsabilidades entre
dominio, acceso a datos, control y presentación.

\section{Propósito y el sentido de «sin framework»}
\label{ap:crud:proposito}

Esta práctica es el complemento aplicado de la Unidad~II. En los
capítulos~\ref{cap:semana5} y~\ref{cap:semana6} se presentaron PHP y
Java en el servidor, y la persistencia con frameworks que generan el
SQL automáticamente. Aquí se hace lo contrario: se escribe el SQL
\textbf{de forma explícita} y se encapsula en una capa DAO, de modo que
el estudiante vea con total claridad qué consulta viaja a la base de
datos y cómo se evita la inyección SQL mediante parámetros enlazados.

\begin{cajadefinicion}{Concepto --- El patrón DAO}
Un \emph{Data Access Object} (DAO) es una clase cuya única
responsabilidad es traducir entre los objetos del dominio (por ejemplo,
\texttt{Estudiante}) y las filas de una tabla. Expone métodos de
negocio (\texttt{listar}, \texttt{buscarPorId}, \texttt{insertar},
\texttt{actualizar}, \texttt{eliminar}) y oculta por completo los
detalles de JDBC, PDO o el cursor de base de datos. Quien usa el DAO no
sabe ---ni necesita saber--- que por dentro hay SQL.
\end{cajadefinicion}

La expresión «sin framework» se refiere de manera estricta a la
\textbf{capa de persistencia}: no se utiliza el ORM (Hibernate/JPA en
Java, el ORM de Django, ni Eloquent en Laravel). En la pista de Django
y en la de Laravel sí se emplea el framework para el enrutamiento y las
plantillas ---es su razón de ser---, pero el acceso a datos se resuelve
con una conexión y consultas propias. Así el mismo patrón DAO se
muestra portable entre lenguajes, y se evidencia que el ORM es una
comodidad, no una obligación.

\begin{cajabuenas}{Las cuatro capas que se repiten en las tres pistas}
\begin{enumerate}
\item \textbf{Dominio:} la clase \texttt{Estudiante} (datos puros).
\item \textbf{Acceso a datos:} una clase de \emph{conexión} y un
\texttt{EstudianteDAO} con SQL parametrizado.
\item \textbf{Control:} un servlet/controlador que recibe la petición,
decide la acción y elige la vista.
\item \textbf{Presentación:} las vistas (JSP, plantilla Django o Blade)
que solo muestran datos, sin lógica de negocio.
\end{enumerate}
\end{cajabuenas}

\section{Modelo de datos común}
\label{ap:crud:modelo}

Las tres aplicaciones operan sobre la misma base de datos PostgreSQL
\texttt{appweb2026ppa} y la misma tabla \texttt{estudiantes}, ya
empleada en la Unidad~II. El esquema se crea una sola vez; las tres
pistas lo reutilizan sin cambios.

{\small\textbf{Crear en (cualquier pista):} \texttt{db/schema.sql}, y
ejecutarlo en PostgreSQL.}
\begin{lstlisting}[style=sql,caption={Esquema compartido de la tabla estudiantes},label={lst:crud-schema}]
CREATE TABLE IF NOT EXISTS estudiantes (
    id      SERIAL       PRIMARY KEY,
    nombre  VARCHAR(100) NOT NULL,
    correo  VARCHAR(120) NOT NULL,
    carrera VARCHAR(80)  NOT NULL,
    creado  TIMESTAMP    NOT NULL DEFAULT NOW()
);
\end{lstlisting}

\begin{cajaadvertencia}{Datos de conexión usados en todo el apéndice}
Servidor \texttt{localhost}, puerto \texttt{5432}, base
\texttt{appweb2026ppa}, usuario \texttt{postgres}, contraseña
\texttt{P@\$7Gr3\$QL}. En un proyecto real, estas credenciales no se
escriben en el código fuente: se leen de variables de entorno o de un
archivo de configuración fuera del control de versiones. Se muestran
aquí en el código solo para que la práctica sea reproducible de
principio a fin.
\end{cajaadvertencia}

% =====================================================================
\section{Pista 1 --- Java con JSP, JSTL y EL (JDBC puro, sin Spring)}
\label{ap:crud:jsp}

Esta pista usa únicamente la plataforma base de Java para la web:
servlets como controladores, JSP con JSTL y EL como vistas, y JDBC para
el acceso a datos. El servidor de aplicaciones es Apache Tomcat~11, que
implementa Jakarta Servlet~6.1, Jakarta Server Pages~4.0 y Jakarta
Expression Language~6.0 \cite{pgjdbc}.

\subsection{Crear el proyecto en IntelliJ IDEA}
\label{ap:crud:jsp:crear}

\begin{enumerate}
\item Menú \texttt{File > New > Project}.
\item En el panel izquierdo, elija \textbf{Jakarta EE} (en IntelliJ
IDEA Ultimate). Si prefiere partir de cero, elija \textbf{Maven
Archetype} con \texttt{maven-archetype-webapp}.
\item \textbf{Name:} \texttt{gestion-estudiantes-jsp}. \textbf{Build
system:} Maven. \textbf{JDK:} 21 (LTS).
\item \textbf{Template:} \emph{Web application}. \textbf{Application
server:} agregue \emph{Tomcat 11} con el botón \texttt{New...} si aún
no está registrado, indicando la carpeta donde descomprimió Tomcat.
\item \texttt{Create}. IntelliJ genera el esqueleto Maven con
empaquetado \texttt{war}.
\end{enumerate}

\begin{cajaejemplo}{Lo que IntelliJ genera automáticamente}
\begin{lstlisting}[style=shell,caption={Esqueleto inicial creado por IntelliJ IDEA},label={lst:jsp-skel}]
gestion-estudiantes-jsp/
|-- pom.xml                         (IntelliJ)
`-- src/
    `-- main/
        |-- java/                   (IntelliJ, vacio)
        |-- resources/              (IntelliJ, vacio)
        `-- webapp/
            |-- WEB-INF/
            |   `-- web.xml         (IntelliJ)
            `-- index.jsp           (IntelliJ)
\end{lstlisting}
\end{cajaejemplo}

\subsection{Estructura final del proyecto}
\label{ap:crud:jsp:estructura}

A partir del esqueleto, el programador crea los paquetes y archivos
marcados como \emph{(programador)}. La organización por capas es la
columna vertebral de toda la práctica.

\begin{lstlisting}[style=shell,caption={Estructura completa de la pista JSP/JSTL/EL},label={lst:jsp-tree}]
gestion-estudiantes-jsp/
|-- pom.xml                                   (programador: dependencias)
|-- db/
|   `-- schema.sql                            (programador)
`-- src/main/
    |-- java/ec/edu/uteq/appweb/
    |   |-- modelo/
    |   |   `-- Estudiante.java               (programador: dominio)
    |   |-- dao/
    |   |   |-- ConexionBD.java               (programador: conexion)
    |   |   `-- EstudianteDAO.java            (programador: acceso a datos)
    |   `-- controlador/
    |       `-- EstudianteServlet.java        (programador: control)
    `-- webapp/
        |-- WEB-INF/
        |   |-- web.xml                       (IntelliJ)
        |   `-- vistas/
        |       |-- lista.jsp                 (programador: vista)
        |       `-- formulario.jsp            (programador: vista)
        `-- index.jsp                         (IntelliJ; redirige al servlet)
\end{lstlisting}

\subsection{Dependencias (pom.xml)}
\label{ap:crud:jsp:pom}

Las API de Servlet y JSP las provee Tomcat en tiempo de ejecución, por
eso su \texttt{scope} es \texttt{provided}. JSTL y el driver de
PostgreSQL sí se empaquetan dentro del WAR.

{\small\textbf{Editar:} \texttt{pom.xml}}
\begin{lstlisting}[style=xml,caption={Dependencias del proyecto JSP},label={lst:jsp-pom}]
<properties>
    <maven.compiler.release>21</maven.compiler.release>
    <project.build.sourceEncoding>UTF-8</project.build.sourceEncoding>
</properties>

<dependencies>
    <!-- API de Jakarta Servlet (la provee Tomcat) -->
    <dependency>
        <groupId>jakarta.servlet</groupId>
        <artifactId>jakarta.servlet-api</artifactId>
        <version>6.1.0</version>
        <scope>provided</scope>
    </dependency>
    <!-- API de Jakarta Server Pages (la provee Tomcat) -->
    <dependency>
        <groupId>jakarta.servlet.jsp</groupId>
        <artifactId>jakarta.servlet.jsp-api</artifactId>
        <version>4.0.0</version>
        <scope>provided</scope>
    </dependency>
    <!-- JSTL: API + implementacion (se empaquetan en el WAR) -->
    <dependency>
        <groupId>jakarta.servlet.jsp.jstl</groupId>
        <artifactId>jakarta.servlet.jsp.jstl-api</artifactId>
        <version>3.0.2</version>
    </dependency>
    <dependency>
        <groupId>org.glassfish.web</groupId>
        <artifactId>jakarta.servlet.jsp.jstl</artifactId>
        <version>3.0.1</version>
    </dependency>
    <!-- Driver JDBC de PostgreSQL -->
    <dependency>
        <groupId>org.postgresql</groupId>
        <artifactId>postgresql</artifactId>
        <version>42.7.11</version>
    </dependency>
</dependencies>
\end{lstlisting}

\subsection{La clase de dominio}
\label{ap:crud:jsp:dominio}

Para que EL pueda leer \texttt{\$\{e.nombre\}} en la JSP, el objeto debe
ser un \emph{JavaBean}: atributos privados y métodos \texttt{get}/%
\texttt{set} públicos. Por eso se usa una clase clásica y no un
\texttt{record}.

{\small\textbf{Crear en:} \texttt{src/main/java/ec/edu/uteq/appweb/modelo/Estudiante.java}}
\begin{lstlisting}[style=java,caption={Estudiante.java (JavaBean del dominio)},label={lst:jsp-modelo}]
package ec.edu.uteq.appweb.modelo;

public class Estudiante {
    private int id;
    private String nombre;
    private String correo;
    private String carrera;

    public Estudiante() { }

    public Estudiante(int id, String nombre, String correo, String carrera) {
        this.id = id;
        this.nombre = nombre;
        this.correo = correo;
        this.carrera = carrera;
    }
    public int getId() { return id; }
    public void setId(int id) { this.id = id; }
    public String getNombre() { return nombre; }
    public void setNombre(String nombre) { this.nombre = nombre; }
    public String getCorreo() { return correo; }
    public void setCorreo(String correo) { this.correo = correo; }
    public String getCarrera() { return carrera; }
    public void setCarrera(String carrera) { this.carrera = carrera; }
}
\end{lstlisting}

\subsection{La conexión a la base de datos}
\label{ap:crud:jsp:conexion}

Desde JDBC~4 el driver se registra solo (vía \emph{ServiceLoader}), de
modo que no hace falta \texttt{Class.forName}. La conexión se centraliza
en una clase con un único método estático.

{\small\textbf{Crear en:} \texttt{src/main/java/ec/edu/uteq/appweb/dao/ConexionBD.java}}
\begin{lstlisting}[style=java,caption={ConexionBD.java (conexion JDBC centralizada)},label={lst:jsp-conexion}]
package ec.edu.uteq.appweb.dao;

import java.sql.Connection;
import java.sql.DriverManager;
import java.sql.SQLException;

public final class ConexionBD {
    private static final String URL =
        "jdbc:postgresql://localhost:5432/appweb2026ppa";
    private static final String USUARIO = "postgres";
    private static final String CLAVE = "P@$7Gr3$QL";

    private ConexionBD() { }

    public static Connection obtener() throws SQLException {
        return DriverManager.getConnection(URL, USUARIO, CLAVE);
    }
}
\end{lstlisting}

\subsection{El DAO con SQL parametrizado}
\label{ap:crud:jsp:dao}

Cada método abre su conexión con \emph{try-with-resources}: así la
conexión, la sentencia preparada (\texttt{Prepared\-Statement}) y el
conjunto de resultados (\texttt{Result\-Set}) se cierran de forma
automática. Todos los datos externos viajan como parámetros
(\texttt{?}), nunca concatenados: así se previene la inyección SQL.

{\small\textbf{Crear en:} \texttt{dao/EstudianteDAO.java} (parte 1 de 2: lectura)}
\begin{lstlisting}[style=java,caption={EstudianteDAO.java --- parte 1: listar y buscar},label={lst:jsp-dao-a}]
package ec.edu.uteq.appweb.dao;

import ec.edu.uteq.appweb.modelo.Estudiante;
import java.sql.Connection;
import java.sql.PreparedStatement;
import java.sql.ResultSet;
import java.sql.SQLException;
import java.util.ArrayList;
import java.util.List;

public class EstudianteDAO {

    public List<Estudiante> listar() throws SQLException {
        String sql = "SELECT id, nombre, correo, carrera "
                   + "FROM estudiantes ORDER BY id";
        List<Estudiante> lista = new ArrayList<>();
        try (Connection cn = ConexionBD.obtener();
             PreparedStatement ps = cn.prepareStatement(sql);
             ResultSet rs = ps.executeQuery()) {
            while (rs.next()) {
                lista.add(mapear(rs));
            }
        }
        return lista;
    }

    public Estudiante buscarPorId(int id) throws SQLException {
        String sql = "SELECT id, nombre, correo, carrera "
                   + "FROM estudiantes WHERE id = ?";
        try (Connection cn = ConexionBD.obtener();
             PreparedStatement ps = cn.prepareStatement(sql)) {
            ps.setInt(1, id);
            try (ResultSet rs = ps.executeQuery()) {
                return rs.next() ? mapear(rs) : null;
            }
        }
    }
\end{lstlisting}

{\small\textbf{Continúa en:} \texttt{dao/EstudianteDAO.java} (parte 2 de 2: escritura)}
\begin{lstlisting}[style=java,caption={EstudianteDAO.java --- parte 2: insertar, actualizar, eliminar},label={lst:jsp-dao-b}]
    public void insertar(Estudiante e) throws SQLException {
        String sql = "INSERT INTO estudiantes (nombre, correo, carrera) "
                   + "VALUES (?, ?, ?)";
        try (Connection cn = ConexionBD.obtener();
             PreparedStatement ps = cn.prepareStatement(sql)) {
            ps.setString(1, e.getNombre());
            ps.setString(2, e.getCorreo());
            ps.setString(3, e.getCarrera());
            ps.executeUpdate();
        }
    }

    public void actualizar(Estudiante e) throws SQLException {
        String sql = "UPDATE estudiantes "
                   + "SET nombre = ?, correo = ?, carrera = ? WHERE id = ?";
        try (Connection cn = ConexionBD.obtener();
             PreparedStatement ps = cn.prepareStatement(sql)) {
            ps.setString(1, e.getNombre());
            ps.setString(2, e.getCorreo());
            ps.setString(3, e.getCarrera());
            ps.setInt(4, e.getId());
            ps.executeUpdate();
        }
    }

    public void eliminar(int id) throws SQLException {
        String sql = "DELETE FROM estudiantes WHERE id = ?";
        try (Connection cn = ConexionBD.obtener();
             PreparedStatement ps = cn.prepareStatement(sql)) {
            ps.setInt(1, id);
            ps.executeUpdate();
        }
    }

    private Estudiante mapear(ResultSet rs) throws SQLException {
        return new Estudiante(
            rs.getInt("id"),
            rs.getString("nombre"),
            rs.getString("correo"),
            rs.getString("carrera"));
    }
}
\end{lstlisting}

\subsection{El controlador (Servlet)}
\label{ap:crud:jsp:servlet}

Un único servlet atiende la ruta \texttt{/estudiantes} y decide la
acción según el parámetro \texttt{accion}. Las lecturas usan
\texttt{doGet}; las escrituras (guardar) usan \texttt{doPost} y
terminan con un \emph{redirect} para evitar reenvíos del formulario.

{\small\textbf{Crear en:} \texttt{.../controlador/EstudianteServlet.java}}
\begin{lstlisting}[style=java,caption={EstudianteServlet.java (control de la peticion)},label={lst:jsp-servlet}]
package ec.edu.uteq.appweb.controlador;

import ec.edu.uteq.appweb.dao.EstudianteDAO;
import ec.edu.uteq.appweb.modelo.Estudiante;
import jakarta.servlet.ServletException;
import jakarta.servlet.annotation.WebServlet;
import jakarta.servlet.http.HttpServlet;
import jakarta.servlet.http.HttpServletRequest;
import jakarta.servlet.http.HttpServletResponse;
import java.io.IOException;
import java.sql.SQLException;
import java.util.List;

@WebServlet("/estudiantes")
public class EstudianteServlet extends HttpServlet {

    private final EstudianteDAO dao = new EstudianteDAO();

    @Override
    protected void doGet(HttpServletRequest req, HttpServletResponse resp)
            throws ServletException, IOException {
        String accion = req.getParameter("accion");
        if (accion == null) { accion = "listar"; }
        try {
            switch (accion) {
                case "nuevo" -> mostrarFormulario(req, resp, null);
                case "editar" -> {
                    int id = Integer.parseInt(req.getParameter("id"));
                    mostrarFormulario(req, resp, dao.buscarPorId(id));
                }
                case "eliminar" -> {
                    dao.eliminar(Integer.parseInt(req.getParameter("id")));
                    resp.sendRedirect(req.getContextPath() + "/estudiantes");
                }
                default -> listar(req, resp);
            }
        } catch (SQLException ex) {
            throw new ServletException("Error de base de datos", ex);
        }
    }

    @Override
    protected void doPost(HttpServletRequest req, HttpServletResponse resp)
            throws ServletException, IOException {
        req.setCharacterEncoding("UTF-8");
        Estudiante e = new Estudiante();
        String id = req.getParameter("id");
        e.setNombre(req.getParameter("nombre"));
        e.setCorreo(req.getParameter("correo"));
        e.setCarrera(req.getParameter("carrera"));
        try {
            if (id == null || id.isBlank()) {
                dao.insertar(e);
            } else {
                e.setId(Integer.parseInt(id));
                dao.actualizar(e);
            }
        } catch (SQLException ex) {
            throw new ServletException("Error al guardar", ex);
        }
        resp.sendRedirect(req.getContextPath() + "/estudiantes");
    }

    private void listar(HttpServletRequest req, HttpServletResponse resp)
            throws SQLException, ServletException, IOException {
        req.setAttribute("lista", dao.listar());
        req.getRequestDispatcher("/WEB-INF/vistas/lista.jsp")
           .forward(req, resp);
    }

    private void mostrarFormulario(HttpServletRequest req,
            HttpServletResponse resp, Estudiante e)
            throws ServletException, IOException {
        req.setAttribute("estudiante", e);
        req.getRequestDispatcher("/WEB-INF/vistas/formulario.jsp")
           .forward(req, resp);
    }
}
\end{lstlisting}

\subsection{Las vistas (JSP + JSTL + EL)}
\label{ap:crud:jsp:vistas}

Las vistas van en \texttt{WEB-INF/vistas/} para que no se puedan abrir
directamente por URL: solo el servlet las despacha. La directiva
\texttt{taglib} con URI \texttt{jakarta.tags.core} habilita JSTL; el
\texttt{c:out} escapa la salida (defensa anti-XSS).

{\small\textbf{Crear en:} \texttt{src/main/webapp/WEB-INF/vistas/lista.jsp}}
\begin{lstlisting}[style=html,caption={lista.jsp (listado con JSTL y EL)},label={lst:jsp-lista}]
<%@ page contentType="text/html; charset=UTF-8" pageEncoding="UTF-8" %>
<%@ taglib prefix="c" uri="jakarta.tags.core" %>
<c:set var="ctx" value="${pageContext.request.contextPath}"/>
<!DOCTYPE html>
<html lang="es">
<head><meta charset="UTF-8"><title>Estudiantes</title></head>
<body>
  <h1>Listado de estudiantes</h1>
  <p><a href="${ctx}/estudiantes?accion=nuevo">Nuevo estudiante</a></p>
  <table border="1">
    <tr><th>ID</th><th>Nombre</th><th>Correo</th>
        <th>Carrera</th><th>Acciones</th></tr>
    <c:forEach var="e" items="${lista}">
      <tr>
        <td>${e.id}</td>
        <td><c:out value="${e.nombre}"/></td>
        <td><c:out value="${e.correo}"/></td>
        <td><c:out value="${e.carrera}"/></td>
        <td>
          <a href="${ctx}/estudiantes?accion=editar&id=${e.id}">Editar</a>
          <a href="${ctx}/estudiantes?accion=eliminar&id=${e.id}"
             onclick="return confirm('Eliminar?');">Eliminar</a>
        </td>
      </tr>
    </c:forEach>
  </table>
</body>
</html>
\end{lstlisting}

{\small\textbf{Crear en:} \texttt{src/main/webapp/WEB-INF/vistas/formulario.jsp}}
\begin{lstlisting}[style=html,caption={formulario.jsp (alta y edicion)},label={lst:jsp-form}]
<%@ page contentType="text/html; charset=UTF-8" pageEncoding="UTF-8" %>
<%@ taglib prefix="c" uri="jakarta.tags.core" %>
<c:set var="ctx" value="${pageContext.request.contextPath}"/>
<!DOCTYPE html>
<html lang="es">
<head><meta charset="UTF-8"><title>Estudiante</title></head>
<body>
  <h1>
    <c:choose>
      <c:when test="${estudiante != null}">Editar</c:when>
      <c:otherwise>Nuevo</c:otherwise>
    </c:choose> estudiante
  </h1>
  <form method="post" action="${ctx}/estudiantes">
    <input type="hidden" name="id" value="${estudiante.id}">
    <p>Nombre:
       <input type="text" name="nombre"
              value="<c:out value='${estudiante.nombre}'/>" required></p>
    <p>Correo:
       <input type="email" name="correo"
              value="<c:out value='${estudiante.correo}'/>" required></p>
    <p>Carrera:
       <input type="text" name="carrera"
              value="<c:out value='${estudiante.carrera}'/>" required></p>
    <button type="submit">Guardar</button>
    <a href="${ctx}/estudiantes">Cancelar</a>
  </form>
</body>
</html>
\end{lstlisting}

Para que la raíz del sitio lleve al listado, \texttt{index.jsp} solo
redirige:

{\small\textbf{Editar:} \texttt{src/main/webapp/index.jsp}}
\begin{lstlisting}[style=html,caption={index.jsp (redireccion al servlet)},label={lst:jsp-index}]
<% response.sendRedirect(request.getContextPath() + "/estudiantes"); %>
\end{lstlisting}

\subsection{Ejecutar y verificar}
\label{ap:crud:jsp:run}

\begin{enumerate}
\item Ejecute \texttt{db/schema.sql} en PostgreSQL (pestaña
\emph{Database} de IntelliJ o \texttt{psql}).
\item \texttt{Run > Edit Configurations > +} \texttt{> Tomcat Server >
Local}. En la pestaña \emph{Deployment} agregue el artefacto
\texttt{gestion-estudiantes-jsp:war exploded}.
\item Pulse \texttt{Run}. IntelliJ despliega en Tomcat y abre el
navegador en la ruta de contexto \texttt{/gestion-estudiantes-jsp/}
(puerto \texttt{8080} por defecto).
\item Verifique el ciclo completo: crear, listar, editar y eliminar un
estudiante. En la consola no deben aparecer trazas de error.
\end{enumerate}

\begin{cajaadvertencia}{Errores típicos en la pista JSP}
\begin{itemize}
\item \texttt{No suitable driver}: falta la dependencia del driver
PostgreSQL en \texttt{pom.xml} o no se recargó Maven.
\item Etiquetas \texttt{<c:forEach>} que aparecen como texto: el URI del
\texttt{taglib} no es \texttt{jakarta.tags.core}, o falta la
implementación de JSTL (segunda dependencia).
\item \texttt{404} al abrir la JSP directamente: es lo esperado; las
vistas están en \texttt{WEB-INF} y solo se acceden vía el servlet.
\end{itemize}
\end{cajaadvertencia}

% =====================================================================
\section{Pista 2 --- Django con SQL parametrizado (sin ORM)}
\label{ap:crud:django}

Django aporta el enrutamiento, las plantillas y la gestion de la
peticion; el acceso a datos, en cambio, se resuelve con un DAO que
ejecuta SQL a mano mediante el cursor de \texttt{django.db.connection}.
No se definen modelos del ORM para \texttt{estudiantes}. Se usa Django
5.2~LTS sobre Python~3.13 \cite{django52docs}.

\subsection{Crear el proyecto en IntelliJ IDEA}
\label{ap:crud:django:crear}

\begin{enumerate}
\item Active el plugin \emph{Python} (\texttt{Settings > Plugins}) en
IntelliJ IDEA Ultimate.
\item \texttt{File > New > Project > Python}. Defina un entorno virtual
(\emph{Virtualenv}) con el interprete Python~3.13.
\item Abra la \emph{Terminal} integrada e instale las dependencias:
\end{enumerate}

\begin{lstlisting}[style=shell,caption={Instalacion de Django y el adaptador de PostgreSQL},label={lst:dj-pip}]
pip install "Django>=5.2,<6.0" "psycopg[binary]"
django-admin startproject gestion .
python manage.py startapp estudiantes
\end{lstlisting}

\begin{cajaejemplo}{Lo que generan startproject y startapp}
\begin{lstlisting}[style=shell,caption={Esqueleto creado por django-admin},label={lst:dj-skel}]
gestion/                          (raiz del proyecto)
|-- manage.py                     (django-admin)
|-- gestion/                      (django-admin: configuracion)
|   |-- __init__.py
|   |-- settings.py
|   |-- urls.py
|   |-- asgi.py
|   `-- wsgi.py
`-- estudiantes/                  (django-admin: la app)
    |-- __init__.py
    |-- admin.py
    |-- apps.py
    |-- migrations/
    |-- models.py                 (NO se usa: no hay ORM)
    |-- tests.py
    `-- views.py
\end{lstlisting}
\end{cajaejemplo}

\subsection{Estructura final del proyecto}
\label{ap:crud:django:estructura}

\begin{lstlisting}[style=shell,caption={Estructura completa de la pista Django},label={lst:dj-tree}]
gestion/
|-- manage.py                                 (django-admin)
|-- db/schema.sql                             (programador)
|-- gestion/
|   |-- settings.py                           (programador: BD + app)
|   `-- urls.py                               (programador: incluye la app)
`-- estudiantes/
    |-- dominio.py                            (programador: dominio)
    |-- dao.py                                (programador: acceso a datos)
    |-- views.py                              (programador: control)
    |-- urls.py                               (programador: rutas de la app)
    `-- templates/estudiantes/
        |-- lista.html                        (programador: vista)
        `-- formulario.html                   (programador: vista)
\end{lstlisting}

\subsection{Configuracion: base de datos y app}
\label{ap:crud:django:settings}

{\small\textbf{Editar:} \texttt{gestion/settings.py}}
\begin{lstlisting}[style=python,caption={settings.py (conexion PostgreSQL y registro de la app)},label={lst:dj-settings}]
INSTALLED_APPS = [
    "django.contrib.staticfiles",
    "estudiantes",
]

DATABASES = {
    "default": {
        "ENGINE": "django.db.backends.postgresql",
        "NAME": "appweb2026ppa",
        "USER": "postgres",
        "PASSWORD": "P@$7Gr3$QL",
        "HOST": "localhost",
        "PORT": "5432",
    }
}
\end{lstlisting}

\subsection{El dominio}
\label{ap:crud:django:dominio}

{\small\textbf{Crear en:} \texttt{estudiantes/dominio.py}}
\begin{lstlisting}[style=python,caption={dominio.py (la clase Estudiante)},label={lst:dj-dominio}]
from dataclasses import dataclass


@dataclass
class Estudiante:
    id: int | None
    nombre: str
    correo: str
    carrera: str
\end{lstlisting}

\subsection{El DAO con SQL parametrizado}
\label{ap:crud:django:dao}

El cursor de Django usa marcadores \texttt{\%s} (estilo DB-API). Al
pasar los valores como segundo argumento de \texttt{execute}, la
biblioteca los enlaza de forma segura: no hay concatenacion, no hay
inyeccion SQL.

{\small\textbf{Crear en:} \texttt{estudiantes/dao.py}}
\begin{lstlisting}[style=python,caption={dao.py (CRUD con cursor, sin el ORM)},label={lst:dj-dao}]
from django.db import connection
from .dominio import Estudiante


class EstudianteDAO:
    """Acceso a datos con SQL parametrizado, sin el ORM de Django."""

    def listar(self):
        sql = ("SELECT id, nombre, correo, carrera "
               "FROM estudiantes ORDER BY id")
        with connection.cursor() as cur:
            cur.execute(sql)
            return [self._mapear(fila) for fila in cur.fetchall()]

    def buscar_por_id(self, id_est):
        sql = ("SELECT id, nombre, correo, carrera "
               "FROM estudiantes WHERE id = %s")
        with connection.cursor() as cur:
            cur.execute(sql, [id_est])
            fila = cur.fetchone()
            return self._mapear(fila) if fila else None

    def insertar(self, e):
        sql = ("INSERT INTO estudiantes (nombre, correo, carrera) "
               "VALUES (%s, %s, %s)")
        with connection.cursor() as cur:
            cur.execute(sql, [e.nombre, e.correo, e.carrera])

    def actualizar(self, e):
        sql = ("UPDATE estudiantes "
               "SET nombre = %s, correo = %s, carrera = %s WHERE id = %s")
        with connection.cursor() as cur:
            cur.execute(sql, [e.nombre, e.correo, e.carrera, e.id])

    def eliminar(self, id_est):
        sql = "DELETE FROM estudiantes WHERE id = %s"
        with connection.cursor() as cur:
            cur.execute(sql, [id_est])

    @staticmethod
    def _mapear(fila):
        return Estudiante(id=fila[0], nombre=fila[1],
                          correo=fila[2], carrera=fila[3])
\end{lstlisting}

\subsection{El control (views) y las rutas}
\label{ap:crud:django:views}

{\small\textbf{Crear en:} \texttt{estudiantes/views.py}}
\begin{lstlisting}[style=python,caption={views.py (control de la peticion)},label={lst:dj-views}]
from django.shortcuts import render, redirect
from .dao import EstudianteDAO
from .dominio import Estudiante

dao = EstudianteDAO()


def listar(request):
    return render(request, "estudiantes/lista.html",
                  {"estudiantes": dao.listar()})


def nuevo(request):
    return render(request, "estudiantes/formulario.html",
                  {"estudiante": None})


def editar(request, id_est):
    return render(request, "estudiantes/formulario.html",
                  {"estudiante": dao.buscar_por_id(id_est)})


def guardar(request):
    id_txt = request.POST.get("id")
    e = Estudiante(
        id=int(id_txt) if id_txt else None,
        nombre=request.POST.get("nombre"),
        correo=request.POST.get("correo"),
        carrera=request.POST.get("carrera"))
    if e.id:
        dao.actualizar(e)
    else:
        dao.insertar(e)
    return redirect("estudiantes:listar")


def eliminar(request, id_est):
    dao.eliminar(id_est)
    return redirect("estudiantes:listar")
\end{lstlisting}

{\small\textbf{Crear en:} \texttt{estudiantes/urls.py}}
\begin{lstlisting}[style=python,caption={urls.py de la app},label={lst:dj-urls-app}]
from django.urls import path
from . import views

app_name = "estudiantes"

urlpatterns = [
    path("", views.listar, name="listar"),
    path("nuevo/", views.nuevo, name="nuevo"),
    path("editar/<int:id_est>/", views.editar, name="editar"),
    path("guardar/", views.guardar, name="guardar"),
    path("eliminar/<int:id_est>/", views.eliminar, name="eliminar"),
]
\end{lstlisting}

{\small\textbf{Editar:} \texttt{gestion/urls.py}}
\begin{lstlisting}[style=python,caption={urls.py del proyecto (incluye la app)},label={lst:dj-urls-root}]
from django.urls import path, include

urlpatterns = [
    path("estudiantes/", include("estudiantes.urls")),
]
\end{lstlisting}

\subsection{Las plantillas}
\label{ap:crud:django:templates}

El formulario incluye \texttt{\{\% csrf\_token \%\}}: Django rechaza los
POST sin ese testigo, lo que protege contra CSRF.

{\small\textbf{Crear en:} \texttt{estudiantes/templates/estudiantes/lista.html}}
\begin{lstlisting}[style=html,caption={lista.html (plantilla de listado)},label={lst:dj-lista}]
<!DOCTYPE html>
<html lang="es">
<head><meta charset="UTF-8"><title>Estudiantes</title></head>
<body>
  <h1>Listado de estudiantes</h1>
  <p><a href="{% url 'estudiantes:nuevo' %}">Nuevo estudiante</a></p>
  <table border="1">
    <tr><th>ID</th><th>Nombre</th><th>Correo</th>
        <th>Carrera</th><th>Acciones</th></tr>
    {% for e in estudiantes %}
      <tr>
        <td>{{ e.id }}</td>
        <td>{{ e.nombre }}</td>
        <td>{{ e.correo }}</td>
        <td>{{ e.carrera }}</td>
        <td>
          <a href="{% url 'estudiantes:editar' e.id %}">Editar</a>
          <a href="{% url 'estudiantes:eliminar' e.id %}"
             onclick="return confirm('Eliminar?');">Eliminar</a>
        </td>
      </tr>
    {% endfor %}
  </table>
</body>
</html>
\end{lstlisting}

{\small\textbf{Crear en:} \texttt{estudiantes/templates/estudiantes/formulario.html}}
\begin{lstlisting}[style=html,caption={formulario.html (alta y edicion)},label={lst:dj-form}]
<!DOCTYPE html>
<html lang="es">
<head><meta charset="UTF-8"><title>Estudiante</title></head>
<body>
  <h1>{% if estudiante %}Editar{% else %}Nuevo{% endif %} estudiante</h1>
  <form method="post" action="{% url 'estudiantes:guardar' %}">
    {% csrf_token %}
    <input type="hidden" name="id"
           value="{% if estudiante %}{{ estudiante.id }}{% endif %}">
    <p>Nombre: <input type="text" name="nombre"
       value="{% if estudiante %}{{ estudiante.nombre }}{% endif %}" required></p>
    <p>Correo: <input type="email" name="correo"
       value="{% if estudiante %}{{ estudiante.correo }}{% endif %}" required></p>
    <p>Carrera: <input type="text" name="carrera"
       value="{% if estudiante %}{{ estudiante.carrera }}{% endif %}" required></p>
    <button type="submit">Guardar</button>
    <a href="{% url 'estudiantes:listar' %}">Cancelar</a>
  </form>
</body>
</html>
\end{lstlisting}

\subsection{Ejecutar y verificar}
\label{ap:crud:django:run}

\begin{enumerate}
\item Ejecute \texttt{db/schema.sql} en PostgreSQL.
\item En la Terminal: \texttt{python manage.py runserver}.
\item Abra \texttt{http://localhost:8000/estudiantes/} y pruebe crear,
listar, editar y eliminar.
\end{enumerate}

\begin{cajaadvertencia}{Errores tipicos en la pista Django}
\begin{itemize}
\item \texttt{ModuleNotFoundError: psycopg}: falta instalar
\texttt{psycopg[binary]} en el entorno virtual activo.
\item \texttt{TemplateDoesNotExist}: la plantilla no esta en
\texttt{estudiantes/templates/estudiantes/} o falta registrar la app en
\texttt{INSTALLED\_APPS}.
\item \texttt{403 (CSRF verification failed)}: el formulario no incluye
\texttt{\{\% csrf\_token \%\}}.
\end{itemize}
\end{cajaadvertencia}

% =====================================================================
\section{Pista 3 --- Laravel con PDO (sin Eloquent)}
\label{ap:crud:laravel}

Laravel aporta el enrutamiento, los controladores y las plantillas
Blade; el acceso a datos se resuelve con un DAO que usa PDO
directamente, sin Eloquent. Se usa Laravel~13 sobre PHP~8.4
\cite{laravel13docs,phppdo}.

\subsection{Crear el proyecto en IntelliJ IDEA}
\label{ap:crud:laravel:crear}

\begin{enumerate}
\item Active el plugin \emph{PHP} en IntelliJ IDEA Ultimate y configure
el interprete PHP~8.4 (\texttt{Settings > PHP}). Asegure la extension
\texttt{pdo\_pgsql} habilitada.
\item Cree el proyecto con Composer desde la Terminal:
\end{enumerate}

\begin{lstlisting}[style=shell,caption={Creacion del proyecto Laravel},label={lst:lv-create}]
composer create-project laravel/laravel gestion-estudiantes
cd gestion-estudiantes
php artisan key:generate
\end{lstlisting}

\begin{cajaejemplo}{Lo que genera Composer (estructura esencial)}
\begin{lstlisting}[style=shell,caption={Esqueleto creado por laravel/laravel},label={lst:lv-skel}]
gestion-estudiantes/
|-- app/
|   |-- Http/Controllers/Controller.php      (Composer)
|   |-- Models/                              (Composer; no se usa Eloquent)
|   `-- Providers/
|-- bootstrap/  config/  database/
|-- public/index.php                         (Composer: punto de entrada)
|-- resources/views/                         (Composer)
|-- routes/web.php                           (Composer)
|-- .env                                     (Composer: configuracion)
|-- composer.json
`-- artisan
\end{lstlisting}
\end{cajaejemplo}

\subsection{Estructura final del proyecto}
\label{ap:crud:laravel:estructura}

Las clases del DAO y del dominio se ubican en carpetas nuevas dentro de
\texttt{app/}, aprovechando el autocargado PSR-4 (\texttt{App\textbackslash}).

\begin{lstlisting}[style=shell,caption={Estructura completa de la pista Laravel},label={lst:lv-tree}]
gestion-estudiantes/
|-- db/schema.sql                             (programador)
|-- app/
|   |-- Dominio/Estudiante.php                (programador: dominio)
|   |-- Dao/
|   |   |-- ConexionBD.php                     (programador: conexion PDO)
|   |   `-- EstudianteDao.php                  (programador: acceso a datos)
|   `-- Http/Controllers/
|       `-- EstudianteController.php           (programador: control)
|-- routes/web.php                             (programador: rutas)
`-- resources/views/estudiantes/
    |-- lista.blade.php                        (programador: vista)
    `-- formulario.blade.php                   (programador: vista)
\end{lstlisting}

\subsection{La conexion (PDO)}
\label{ap:crud:laravel:conexion}

{\small\textbf{Crear en:} \texttt{app/Dao/ConexionBD.php}}
\begin{lstlisting}[style=php,caption={ConexionBD.php (conexion PDO reutilizable)},label={lst:lv-conexion}]
<?php

namespace App\Dao;

use PDO;

class ConexionBD
{
    private static ?PDO $pdo = null;

    public static function obtener(): PDO
    {
        if (self::$pdo === null) {
            $dsn = 'pgsql:host=localhost;port=5432;dbname=appweb2026ppa';
            self::$pdo = new PDO($dsn, 'postgres', 'P@$7Gr3$QL', [
                PDO::ATTR_ERRMODE => PDO::ERRMODE_EXCEPTION,
                PDO::ATTR_DEFAULT_FETCH_MODE => PDO::FETCH_ASSOC,
            ]);
        }
        return self::$pdo;
    }
}
\end{lstlisting}

\subsection{El dominio}
\label{ap:crud:laravel:dominio}

{\small\textbf{Crear en:} \texttt{app/Dominio/Estudiante.php}}
\begin{lstlisting}[style=php,caption={Estudiante.php (dominio con propiedades promovidas)},label={lst:lv-dominio}]
<?php

namespace App\Dominio;

class Estudiante
{
    public function __construct(
        public ?int $id,
        public string $nombre,
        public string $correo,
        public string $carrera
    ) {}
}
\end{lstlisting}

\subsection{El DAO con SQL parametrizado}
\label{ap:crud:laravel:dao}

PDO usa marcadores con nombre (\texttt{:nombre}). El metodo
\texttt{execute} recibe el arreglo de valores y los enlaza de forma
segura.

{\small\textbf{Crear en:} \texttt{app/Dao/EstudianteDao.php}}
\begin{lstlisting}[style=php,caption={EstudianteDao.php (CRUD con PDO, sin Eloquent)},label={lst:lv-dao}]
<?php

namespace App\Dao;

use App\Dominio\Estudiante;

class EstudianteDao
{
    public function listar(): array
    {
        $sql = 'SELECT id, nombre, correo, carrera '
             . 'FROM estudiantes ORDER BY id';
        $stmt = ConexionBD::obtener()->query($sql);
        return array_map([$this, 'mapear'], $stmt->fetchAll());
    }

    public function buscarPorId(int $id): ?Estudiante
    {
        $sql = 'SELECT id, nombre, correo, carrera '
             . 'FROM estudiantes WHERE id = :id';
        $stmt = ConexionBD::obtener()->prepare($sql);
        $stmt->execute([':id' => $id]);
        $fila = $stmt->fetch();
        return $fila ? $this->mapear($fila) : null;
    }

    public function insertar(Estudiante $e): void
    {
        $sql = 'INSERT INTO estudiantes (nombre, correo, carrera) '
             . 'VALUES (:nombre, :correo, :carrera)';
        $stmt = ConexionBD::obtener()->prepare($sql);
        $stmt->execute([
            ':nombre' => $e->nombre,
            ':correo' => $e->correo,
            ':carrera' => $e->carrera,
        ]);
    }

    public function actualizar(Estudiante $e): void
    {
        $sql = 'UPDATE estudiantes SET nombre = :nombre, '
             . 'correo = :correo, carrera = :carrera WHERE id = :id';
        $stmt = ConexionBD::obtener()->prepare($sql);
        $stmt->execute([
            ':nombre' => $e->nombre,
            ':correo' => $e->correo,
            ':carrera' => $e->carrera,
            ':id' => $e->id,
        ]);
    }

    public function eliminar(int $id): void
    {
        $sql = 'DELETE FROM estudiantes WHERE id = :id';
        $stmt = ConexionBD::obtener()->prepare($sql);
        $stmt->execute([':id' => $id]);
    }

    private function mapear(array $fila): Estudiante
    {
        return new Estudiante(
            (int) $fila['id'],
            $fila['nombre'],
            $fila['correo'],
            $fila['carrera']
        );
    }
}
\end{lstlisting}

\subsection{El control y las rutas}
\label{ap:crud:laravel:control}

{\small\textbf{Crear en:} \texttt{app/Http/Controllers/EstudianteController.php}}
\begin{lstlisting}[style=php,caption={EstudianteController.php (control de la peticion)},label={lst:lv-controller}]
<?php

namespace App\Http\Controllers;

use App\Dao\EstudianteDao;
use App\Dominio\Estudiante;
use Illuminate\Http\Request;

class EstudianteController extends Controller
{
    private EstudianteDao $dao;

    public function __construct()
    {
        $this->dao = new EstudianteDao();
    }

    public function index()
    {
        return view('estudiantes.lista',
            ['estudiantes' => $this->dao->listar()]);
    }

    public function nuevo()
    {
        return view('estudiantes.formulario', ['estudiante' => null]);
    }

    public function editar(int $id)
    {
        return view('estudiantes.formulario',
            ['estudiante' => $this->dao->buscarPorId($id)]);
    }

    public function guardar(Request $request)
    {
        $datos = $request->validate([
            'id' => 'nullable|integer',
            'nombre' => 'required|string|max:100',
            'correo' => 'required|email|max:120',
            'carrera' => 'required|string|max:80',
        ]);
        $e = new Estudiante(
            $datos['id'] ?? null,
            $datos['nombre'], $datos['correo'], $datos['carrera']);
        $e->id ? $this->dao->actualizar($e) : $this->dao->insertar($e);
        return redirect()->route('estudiantes.index');
    }

    public function eliminar(int $id)
    {
        $this->dao->eliminar($id);
        return redirect()->route('estudiantes.index');
    }
}
\end{lstlisting}

{\small\textbf{Editar:} \texttt{routes/web.php}}
\begin{lstlisting}[style=php,caption={web.php (rutas del CRUD)},label={lst:lv-routes}]
<?php

use App\Http\Controllers\EstudianteController;
use Illuminate\Support\Facades\Route;

Route::get('/estudiantes', [EstudianteController::class, 'index'])
    ->name('estudiantes.index');
Route::get('/estudiantes/nuevo', [EstudianteController::class, 'nuevo'])
    ->name('estudiantes.nuevo');
Route::get('/estudiantes/{id}/editar', [EstudianteController::class, 'editar'])
    ->name('estudiantes.editar');
Route::post('/estudiantes/guardar', [EstudianteController::class, 'guardar'])
    ->name('estudiantes.guardar');
Route::post('/estudiantes/{id}/eliminar', [EstudianteController::class, 'eliminar'])
    ->name('estudiantes.eliminar');
\end{lstlisting}

\subsection{Las vistas Blade}
\label{ap:crud:laravel:vistas}

La directiva \texttt{@csrf} inserta el testigo anti-CSRF; las llaves
dobles \texttt{\{\{ \}\}} escapan la salida (anti-XSS).

{\small\textbf{Crear en:} \texttt{resources/views/estudiantes/lista.blade.php}}
\begin{lstlisting}[style=html,caption={lista.blade.php (listado)},label={lst:lv-lista}]
<!DOCTYPE html>
<html lang="es">
<head><meta charset="UTF-8"><title>Estudiantes</title></head>
<body>
  <h1>Listado de estudiantes</h1>
  <p><a href="{{ route('estudiantes.nuevo') }}">Nuevo estudiante</a></p>
  <table border="1">
    <tr><th>ID</th><th>Nombre</th><th>Correo</th>
        <th>Carrera</th><th>Acciones</th></tr>
    @foreach ($estudiantes as $e)
      <tr>
        <td>{{ $e->id }}</td>
        <td>{{ $e->nombre }}</td>
        <td>{{ $e->correo }}</td>
        <td>{{ $e->carrera }}</td>
        <td>
          <a href="{{ route('estudiantes.editar', $e->id) }}">Editar</a>
          <form method="post" style="display:inline"
                action="{{ route('estudiantes.eliminar', $e->id) }}"
                onsubmit="return confirm('Eliminar?');">
            @csrf
            <button type="submit">Eliminar</button>
          </form>
        </td>
      </tr>
    @endforeach
  </table>
</body>
</html>
\end{lstlisting}

{\small\textbf{Crear en:} \texttt{resources/views/estudiantes/formulario.blade.php}}
\begin{lstlisting}[style=html,caption={formulario.blade.php (alta y edicion)},label={lst:lv-form}]
<!DOCTYPE html>
<html lang="es">
<head><meta charset="UTF-8"><title>Estudiante</title></head>
<body>
  <h1>{{ $estudiante ? 'Editar' : 'Nuevo' }} estudiante</h1>
  <form method="post" action="{{ route('estudiantes.guardar') }}">
    @csrf
    <input type="hidden" name="id" value="{{ $estudiante->id ?? '' }}">
    <p>Nombre: <input type="text" name="nombre"
       value="{{ $estudiante->nombre ?? '' }}" required></p>
    <p>Correo: <input type="email" name="correo"
       value="{{ $estudiante->correo ?? '' }}" required></p>
    <p>Carrera: <input type="text" name="carrera"
       value="{{ $estudiante->carrera ?? '' }}" required></p>
    <button type="submit">Guardar</button>
    <a href="{{ route('estudiantes.index') }}">Cancelar</a>
  </form>
</body>
</html>
\end{lstlisting}

\subsection{Ejecutar y verificar}
\label{ap:crud:laravel:run}

\begin{enumerate}
\item Ejecute \texttt{db/schema.sql} en PostgreSQL.
\item En la Terminal: \texttt{php artisan serve}.
\item Abra \texttt{http://localhost:8000/estudiantes} y pruebe el ciclo
CRUD completo.
\end{enumerate}

\begin{cajaadvertencia}{Errores tipicos en la pista Laravel}
\begin{itemize}
\item \texttt{could not find driver}: falta habilitar la extension
\texttt{pdo\_pgsql} en \texttt{php.ini}.
\item \texttt{419 (Page Expired)}: el formulario no incluye
\texttt{@csrf}.
\item \texttt{Class not found}: el espacio de nombres del archivo no
coincide con su carpeta (PSR-4); revise \texttt{namespace} y la ruta.
\end{itemize}
\end{cajaadvertencia}

% =====================================================================
\section{Los tres enfoques comparados}
\label{ap:crud:comparacion}

Las tres pistas resuelven el mismo problema con la misma arquitectura
en cuatro capas; lo que cambia es la sintaxis y las piezas que aporta
cada plataforma. La tabla~\ref{tab:crud-comparacion} resume las
correspondencias.

\begin{table}[H]
\centering
\caption{Correspondencia entre las tres pistas del CRUD con DAO}
\label{tab:crud-comparacion}
\small
\renewcommand{\arraystretch}{1.3}
\begin{tabular}{>{\raggedright\arraybackslash}p{2.8cm}
                >{\raggedright\arraybackslash}p{3.4cm}
                >{\raggedright\arraybackslash}p{3.2cm}
                >{\raggedright\arraybackslash}p{3.2cm}}
\toprule
\textbf{Aspecto} & \textbf{JSP/JSTL/EL} & \textbf{Django} & \textbf{Laravel} \\
\midrule
Lenguaje / versión & Java 21 & Python 3.13 & PHP 8.4 \\
Servidor de ejecución & Tomcat 11 & runserver (dev) & artisan serve (dev) \\
Conexión a la BD & JDBC \texttt{DriverManager} & cursor de \texttt{connection} & PDO \\
Marcador de parámetro & \texttt{?} & \texttt{\%s} & \texttt{:nombre} \\
Control & Servlet & función \emph{view} & Controller \\
Presentación & JSP + JSTL + EL & plantilla Django & Blade \\
Anti-XSS en la vista & \texttt{c:out} & \texttt{\{\{ \}\}} & \texttt{\{\{ \}\}} \\
Anti-CSRF & (manual) & \texttt{csrf\_token} & \texttt{@csrf} \\
\bottomrule
\end{tabular}
\end{table}

\begin{cajabuenas}{Lo que el estudiante debe llevarse de esta práctica}
\begin{enumerate}
\item El patrón DAO es independiente del lenguaje: una clase de
conexión y un DAO con métodos de negocio resuelven la persistencia en
cualquier pila.
\item El SQL parametrizado (\texttt{?}, \texttt{\%s} o \texttt{:nombre})
es la defensa primaria contra la inyección SQL en las tres pilas.
\item El controlador nunca contiene SQL ni la vista contiene lógica de
negocio: esa separación es lo que hace mantenible la aplicación.
\item El ORM ahorra código, pero comprender el SQL subyacente ---como
en este apéndice--- es lo que permite depurar y optimizar cuando el ORM
no basta.
\end{enumerate}
\end{cajabuenas}

\section{Lista de verificación final}
\label{ap:crud:checklist}

\begin{enumerate}
\item La tabla \texttt{estudiantes} existe en \texttt{appweb2026ppa}.
\item Las cuatro capas (dominio, DAO, control, vista) están separadas
en archivos distintos.
\item Todas las consultas usan parámetros enlazados; no hay
concatenación de cadenas en el SQL.
\item Las vistas escapan la salida del usuario.
\item El ciclo CRUD completo (crear, listar, editar, eliminar) funciona
de extremo a extremo en las tres pistas.
\end{enumerate}
